\chapter{Introducción} \label{ch:introduccion}

% ============================================================
%  BLOQUE A — RELEVANCIA DEL FENÓMENO
%  Propósito: Establecer por qué los ciclones extratropicales
%  del Atlántico Sur importan a escala global y regional.
%  Justifica la elección del objeto de estudio antes de
%  cualquier revisión bibliográfica.
% ============================================================

% --- A1: Rol climático y dinámico global ---
Desde una perspectiva dinámica global, el análisis de los \emph{storm tracks} realizado por \citet{hoskins2005new} muestra a los ciclones extratropicales como el mecanismo dominante para el transporte latitudinal de energía en el Hemisferio Sur, equilibrando el gradiente térmico planetario. La cuantificación de estos transportes ha sido posible gracias al desarrollo de algoritmos de seguimiento objetivo, donde \citet{sinclair1994objective} estableció el precedente fundamental al utilizar la vorticidad relativa para identificar y rastrear sistemas independientemente del flujo de fondo.

% --- A2: Impacto regional y relevancia societal ---
En el hemisferio sur, estos sistemas representan el fenómeno de severidad geofísica dominante para la zona costera regional \citep{gramcianinov2023impact}. En el litoral sur de Brasil, esta amenaza se traduce en riesgos meteoceanográficos severos y pérdidas humanas documentadas: \citet{bartolomei2024extremos} reportan la letalidad causada por ciclones extratropicales invernales en esa región, mientras que \citet{miranda2026patterns} confirman que estos sistemas son los principales motores de tormentas costeras en el Atlántico Sur. Su influencia se extiende además al estado del mar: \citet{sasaki2021intraseasonal} demuestran que estos ciclones modulan la variabilidad intraestacional de las olas en el Atlántico Sur occidental, donde la configuración del viento en superficie determina la transferencia de energía hacia el oleaje, amplificando el riesgo costero más allá de la duración del evento sinóptico. El viento en superficie y la precipitación son, en conjunto, las dos variables más empleadas para caracterizar los impactos de estos sistemas en la costa, por lo que abordarlas de manera conjunta resulta natural para este trabajo, en línea con otros enfoques que analizan más de una variable de superficie simultáneamente \citep{zscheischler2018future}.

% ============================================================
%  BLOQUE B — ESTADO DEL ARTE: LO QUE SE SABE
%  Propósito: Sintetizar el conocimiento consolidado en cuatro
%  dimensiones: (B1) climatología de trayectorias, (B2)
%  diversidad estructural, (B3) forzante orográfico y
%  regionalización, y (B4) papel de la humedad y los
%  transportes diabáticos. Este bloque construye el andamiaje
%  conceptual necesario para delimitar el gap en el Bloque C.
% ============================================================

% --- B1: Climatología de trayectorias y ciclogénesis ---
En el contexto regional, \citet{reboita2010south} proporcionaron una de las primeras climatologías para el Atlántico Sur, identificando tres centros preferenciales de ciclogénesis vinculados a la interacción entre el flujo perturbado por los Andes y los gradientes térmicos oceánicos. Investigaciones posteriores como \citet{gramcianinov2020analysis} e investigaciones recientes como \citet{padilhareinke2026characterization} han caracterizado con robustez la distribución espacio-temporal de estos sistemas, consolidando una base climatológica sobre la frecuencia y localización de la actividad ciclónica en la región.

% --- B2: Diversidad estructural (tipología y ciclo de vida) ---
No obstante, la naturaleza de estas perturbaciones en la región es heterogénea. \citet{hart2003cyclone} demostró que los sistemas transitan a lo largo de un continuo estructural en lugar de pertenecer a categorías estáticas: en el Atlántico Sur esto se expresa, por un lado, en la prevalencia de ciclones subtropicales con estructuras térmicas híbridas documentada por \citet{evans2012climatology}, y por otro, en sistemas de tipo extratropical que siguen el modelo estructural de Shapiro-Keyser, caracterizados por una marcada fractura del frente frío y el desarrollo de una seclusión cálida durante su etapa de madurez \citep{gozzo2013air}. Estas variaciones morfológicas, respaldadas por climatologías regionales de largo plazo \citep{gozzo2014subtropical}, evidencian que el Atlántico Sur se aparta frecuentemente de la conceptualización de ciclones convencionales.

% --- B3: Forzante orográfico y regionalización ---
Comprender la génesis de estos sistemas exige considerar el rol del continente sudamericano como perturbación orográfica. La Cordillera de los Andes interactúa mecánicamente con el flujo de los oestes, forzando la compresión de la columna de vorticidad a barlovento y su estiramiento vertical a sotavento, un mecanismo físico descrito por \citet{gan1994influence} que favorece la formación de centros de baja presión al este de la cordillera. Este forzante topográfico, combinado con la baroclinicidad costera y los gradientes de temperatura superficial del mar, consolida focos de actividad ciclogenética en la costa este de Sudamérica \citep{reboita2026meteorology}.

\label{sec:definicion_regiones}%
Adoptando la regionalización propuesta por \citet{gramcianinov2019properties}, este estudio se centra en tres regiones clave: la región Sur-Sudeste de Brasil (\textit{South Brazil}, SBR), la cuenca de descarga del Río de la Plata (\textit{La Plata Basin}, LPB, por sus siglas en inglés) y la costa sureste de Argentina (\textit{Argentina}, ARG). Estas regiones de ciclogenesis han sido ampliamente documentas por diversos autores como \citet{reboita2010south, crespo2020potential, gramcianinov2019properties, reboita2026meteorology}: (1) la región SBR, favorecida por el efecto de sotavento de los Andes y el chorro de bajos niveles (SALLJ), donde \citet{reboita2022from} documentan frecuentes sistemas híbridos y transiciones tropicales; (2) la región LPB, caracterizada por ciclogénesis explosivas e influenciada por el transporte de humedad del SALLJ, resultando determinante para la inestabilización de la atmósfera \citep{crespo2020potential}; y (3) la región ARG, dominada por la baroclinia del frente polar, donde la dinámica del chorro polar y la propagación de ondas de Rossby constituyen los principales mecanismos de desarrollo \citep{simmonds2000variability}. Esta estratificación responde a regímenes dinámicos distintos, mientras que los sistemas en SBR y LPB exhiben una sensibilidad proyectada hacia la intensificación de los flujos de humedad y la advección cálida en capas bajas, las tendencias en la estructura media de los controladores físicos para la región ARG resultan relativamente débiles \citep{gramcianinov2024early}.

% --- B4: Ingredientes dinámicos: humedad, SALLJ y diabática ---
La variabilidad de la actividad ciclónica es modulada por la corriente en chorro y su interacción con las capas bajas. 
Al este de los Andes, la convergencia de viento en niveles bajos es el principal factor de intensificación de estos sistemas \citep{vera2002cold}, favoreciendo las condiciones que provocan precipitaciones intensas en el sureste de Sudamérica. 
Simultáneamente, el transporte meridional de humedad desde los trópicos es mediado por el Chorro de Bajos Niveles de Sudamérica (SALLJ). 
Como señalan \citet{reboita2022from}, la advección de aire cálido y húmedo en el flanco oriental del ciclón no solo inestabiliza la atmósfera, sino que actúa como una fuente diabática que, al liberar calor latente, refuerza la intensidad del sistema.

En el Hemisferio Norte (HN), el flujo ascendente del Cinturón Transportador Cálido (WCB, por sus siglas en inglés) se describe típicamente hacia el noreste \citep{blanchard2021warm}; en el Hemisferio Sur (HS) este flujo se orienta preferencialmente hacia el sureste del sistema.
% de modo que los patrones estructurales del HN deben interpretarse como una imagen especular latitudinal al trasladarse al HS.
El WCB no solo produce precipitación a gran escala, sino que frecuentemente alberga núcleos de convección intensa embebida, generando tasas de precipitación extrema \citep{oertel2021warm}. Investigaciones recientes de \citet{gentile2025response} anticipan, mediante modelaciones de alta resolución, que el sector cálido de los ciclones extratropicales más intensos se intensificará, incrementando significativamente los vientos y la precipitación extrema durante sus etapas de desarrollo.

% ============================================================
%  BLOQUE C — GAP DE CONOCIMIENTO: LO QUE NO SE SABE
%  Propósito: Articular de forma unificada y explícita la
%  brecha científica que motiva esta investigación. El gap
%  se construye en tres dimensiones acumulativas:
%  (C1) insuficiencia de las métricas escalares de intensidad,
%  (C2) ausencia de caracterización estructural por fase
%  evolutiva, y (C3) limitaciones de resolución de datos.
%  Cada dimensión prepara la justificativa del Bloque D.
% ============================================================

% --- C1: Insuficiencia de las métricas escalares de intensidad ---
Aunque existen múltiples estudios de estos sistemas, persiste la limitación conceptual en la forma en que se cuantifica la severidad de estos sistemas. Como argumenta \citet{corner2025classification} para el Atlántico Norte, la intensidad de un ciclón no puede capturarse con una única métrica, porque la relación entre la intensidad del sistema y los vientos en superficie no es lineal. \citet{sinclair2023relationship} advierten que un ciclón puede ser dinámicamente intenso pero generar poca lluvia si carece del suministro de humedad adecuado, subrayando que los máximos alcanzados no necesariamente coinciden con el mínimo de presión ocurrido. \citet{cardoso2022synoptic} evidencian que los máximos de viento no se distribuyen uniformemente alrededor del centro, sino que se concentran en cuadrantes específicos. En el HN, \citet{eisenstein2023identification} destacan además que los mayores impactos en superficie suelen estar asociados a estructuras de mesoescala específicas, como los cinturones de transporte frío (CCB), cuya localización espacial respecto al centro de vorticidad varía durante el ciclo de vida del sistema. Esta insuficiencia de las métricas escalares se confirma también en el Hemisferio Sur: \citet{pepler2020dimensional} muestran para el sudeste de Australia que los ciclones profundos presentan vientos e intensidad de lluvia sistemáticamente mayores que los ciclones superficiales de magnitud escalar comparable (11.7 frente a 9.6~m\,s$^{-1}$ de viento medio, y más del doble de precipitación acumulada).

% --- C2: Ausencia de caracterización estructural por fase evolutiva ---
Existe un vacío en la comprensión de cómo se estructura espacialmente la distribución de viento y precipitación durante la evolución del ciclón, y cómo varía en las distintas fases del ciclo de vida. 
En Sudamérica, trabajos de \citet{gozzo2013air} y \citet{reboita2022from} han revelado que los ciclones exhiben evoluciones que se apartan de los modelos tradicionales.
La coocurrencia de los máximos de viento y precipitación es una consecuencia estructural del acoplamiento dinámico entre los mecanismos que gobiernan cada campo \citep{chen2025characteristics}, por lo que caracterizar la morfología superficial del sistema exige analizar ambos máximos de forma conjunta. 
No obstante, su estudio exige una descripción de la organización espacial independiente de cada campo y de su covariabilidad estructural a lo largo del ciclo de vida. Recientemente, \citet{han2025system} enfatiza que se requieren marcos de clasificación que vinculen la evolución del sistema con sus manifestaciones extremas de viento y precipitación.

% --- C3: Limitaciones de resolución de datos ---
Esta caracterización estructural ha sido limitada por la resolución espacial de los datos disponibles. \citet{priestley2022improved} demuestran que los modelos subestiman los vientos extremos asociados a ciclones al no resolver explícitamente sus estructuras de mesoescala. Asimismo, \citet{neu2013intercomparison} señalan que la variabilidad en la forma y tamaño de los ciclones extratropicales genera incertidumbres en la determinación de su intensidad en superficie, especialmente en regiones con topografía compleja.

% ============================================================
%  BLOQUE D — JUSTIFICATIVA Y OPORTUNIDAD METODOLÓGICA
%  Propósito: Demostrar que el gap identificado en el Bloque C
%  es ahora abordable gracias a la convergencia de tres
%  avances recientes: (D1) el reanálisis ERA5, (D2) el
%  algoritmo de clasificación de fases CycloPhaser, y
%  (D3) el seguimiento por vorticidad relativa. Este bloque
%  transforma el problema en una investigación viable.
% ============================================================

Este estudio utiliza la base de datos de trayectorias de \citet{coutodesouza2024new}, fundamentada en el seguimiento por vorticidad relativa a 850 hPa ($\zeta_{850}$), variable que permite aislar la circulación del sistema del flujo de fondo y capturar sistemas en etapas tempranas que los algoritmos basados en presión mínima tienden a omitir \citep{sinclair1994objective,sinclair1997objective}. 
Los fundamentos dinámicos y la validación de este criterio para el Hemisferio Sur se desarrollan en la Sección~\ref{sec:tb_vorticidad}. El uso de $\zeta_{850}$ resulta particularmente crítico para esta investigación porque garantiza la representación de la fase incipiente, condición necesaria para una caracterización estructural completa del ciclo de vida.

Para caracterizar la evolución física de la estructura interna, se adopta la segmentación objetiva desarrollada por \citet{coutodesouza2024new} mediante el algoritmo \textit{CycloPhaser}, que define cuatro fases dinámicas discretas, incipiente, intensificación, madurez y decaimiento \citep{deSouza2025CycloPhaser}. 
Este marco supera las limitaciones de las clasificaciones clásicas en el contexto del Atlántico Sur (discutidas en la Sección~\ref{subsec:tb_fases}), permitiendo analizar cómo se transforma la distribución espacial de los campos de superficie a través de las distintas etapas del sistema. Esta estrategia se alinea con propuestas globales recientes como \textit{SyCLoPS} de \citet{han2025system}. 
La caracterización de esta evolución estructural se apoya además en el modelo de Fractura Frontal \citep{shapiro1990life,shapiro1999bridge} como marco de referencia para interpretar la distribución espacial de los máximos de viento y precipitación; su desarrollo se presenta en la Sección~\ref{subsec:tb_modelos}.

La viabilidad de esta caracterización estructural reside en la disponibilidad del reanálisis ERA5 \citep{hersbach2020era5}, cuya resolución espacial ($\sim$31 km) y temporal (horaria) permite resolver los gradientes de presión y los flujos de humedad que generaciones anteriores de datos suavizaban. Como demuestra \citet{priestley2022improved}, una mayor resolución espacial permite capturar de forma más explícita las estructuras de mesoescala del viento que los productos de menor resolución subestiman sistemáticamente. Validaciones específicas para el Atlántico Sur realizadas por \citet{gramcianinov2020analysis} confirman que ERA5 ofrece una representación superior de la intensidad de los vientos superficiales y la estructura de los ciclones en comparación con productos de generaciones previas. No obstante, ERA5 presenta sesgos documentados en la representación de los extremos más intensos, cuya magnitud y consecuencias para la interpretación de los resultados se discuten en la Sección~\ref{sec:datos_era5}.

% ============================================================
%  BLOQUE E — PREGUNTA DE INVESTIGACIÓN, HIPÓTESIS Y
%             ESTRATEGIA ANALÍTICA
%  Propósito: Formular el problema central y la hipótesis
%  de trabajo, e introducir inmediatamente la estrategia
%  estadística (MEOF) que hace la hipótesis testeable.
%  La contigüidad hipótesis–método es clave para la
%  coherencia deductiva exigida en revistas Q1.
% ============================================================

Ante esta complejidad estructural, se plantea la interrogante central de esta investigación: ¿cómo evoluciona espacialmente la distribución de viento y precipitación en los ciclones extratropicales del Atlántico Sur, y qué configuraciones caracterizan esta organización durante las fases incipiente, intensificación, madurez y decaimiento? Se postula que los campos de viento y precipitación presentan patrones de asimetría sistemática que varían entre fases evolutivas, entre regiones de génesis y entre estratos de intensidad dinámica —contrastando la población general de ciclones frente al subconjunto de mayor vorticidad (percentil 90)—, y que dichos patrones no son capturables mediante métricas escalares de intensidad central.

\label{subsec:intro_estrategia_estadistica}
Para testear esta hipótesis, se requieren técnicas estadísticas que trasciendan el análisis univariado y capturen la covariabilidad entre campos físicos distintos. 
Dado que la estructura ciclónica presenta una variabilidad espacial continua donde la circulación del viento y los núcleos de precipitación pueden covaríar durante la evolución del sistema, resulta necesario recurrir a herramientas de reducción dimensional que preserven la interpretabilidad física. En este sentido, \citet{hannachi2007empirical} establece el análisis de Funciones Ortogonales Empíricas Multivariadas (MEOF) como el marco apropiado; a diferencia del análisis univariado de EOF aplicado separadamente a cada campo, la técnica MEOF permite extraer modos de variabilidad inherentemente acoplados entre variables físicamente distintas \citep{liang2018multivariate}, capturando así la organización espacial conjunta de ambos campos. La robustez estadística de los patrones espaciales identificados mediante EOF y MEOF se evalúa además con técnicas de remuestreo (\textit{bootstrap}), lo que permite distinguir señales físicas consistentes del ruido de muestreo propio de cada subconjunto.

La combinación de los datos (ERA5), la detección dinámica ($\zeta_{850}$) y la clasificación objetiva de fases (\textit{CycloPhaser}) ya disponibles en la base de \citet{coutodesouza2024new}, junto con el análisis multivariado (MEOF) desarrollado en este trabajo, constituye la base técnica para cuantificar la evolución de la estructura espacial de los campos de viento y precipitación en la región.
La aplicación de este marco sobre las fases previamente clasificadas permitirá construir prototipos estructurales, sintetizando la morfología típica del ciclón para cada región de génesis y fase evolutiva. 

% ============================================================
%  BLOQUE F — OBJETIVOS
%  Propósito: Traducir la pregunta e hipótesis del Bloque E
%  en metas operacionales concretas y verificables, con
%  correspondencia directa a las secciones metodológicas
%  del manuscrito.
% ============================================================

\section{Objetivos del Proyecto de Investigación}

\subsection{Objetivo General}
Caracterizar la evolución de la estructura espacial del viento a 10 metros y precipitación durante las fases del ciclo de vida de ciclones extratropicales en el Atlántico Sur, caracterizando sus patrones de variabilidad y estimación de densidad probabilística de máximos, con estratificación por región de génesis.

\subsection{Objetivos Específicos}
\begin{enumerate}
    \item Consolidar una base de datos lagrangiana que asocie los campos horarios de alta resolución de viento y precipitación (ERA5) con las trayectorias y fases dinámicas (incipiente, intensificación, madurez y decaimiento) ya clasificadas por \citet{coutodesouza2024new} para los ciclones extratropicales del Atlántico Sur, filtrando los sistemas con ciclo de vida completo y estratificando la muestra resultante por región de génesis (SBR, LPB y ARG) y por intensidad dinámica.
    
    \item Cuantificar la distribución estadística y la asimetría espacial de los máximos de viento a 10 m y precipitación mediante estimación de densidad por núcleos, caracterizando tanto la distribución de probabilidad de sus magnitudes (PDF unidimensional) como la localización preferencial de dichos máximos respecto al centro del ciclón (PDF espacial).
    
    \item Identificar los modos dominantes de variabilidad espacial de los campos de viento y precipitación mediante Funciones Ortogonales Empíricas univariadas (EOF) y multivariadas (MEOF), evaluando el patrón estructural individual de cada variable y la covariabilidad acoplada.
    
    \item Construir prototipos estructurales sintéticos por fase evolutiva y región de génesis, integrando los campos medios condicionados al modo dominante de covariabilidad (MEOF) con las densidades probabilísticas de localización de máximos (PDFe), para cuantificar la magnitud característica de viento y precipitación.
\end{enumerate}